%==============================================================================
%  CAPÍTULO II
%==============================================================================
\chapter{El nuevo estatuto orgánico de los deudores tras la Ley \Nro 21.563}
\label{cap:estatuto-deudores}

\fichaCapitulo{Redefinición subjetiva del concurso.}{Categorías de Empresa y Persona
Deudora, y umbrales de tamaño de la Ley \Nro 20.416.}

%==============================================================================
\bloqueTeorico
%==============================================================================

\subsection{El tamaño de la unidad económica como criterio de asignación procesal}
\pendiente{Desarrollar: por qué el legislador abandona un procedimiento único y adopta
la diferenciación por tamaño e ingresos.}

\subsection{La asimetría de costos de transacción bajo el texto original de 2014}
\pendiente{Desarrollar: liquidación sistemática e ineficiente de MIPEs viables por
inaccesibilidad económica del procedimiento ordinario.}

\subsection{El estatuto diferenciado como técnica de reducción de barreras de acceso}
\pendiente{Desarrollar.}

%==============================================================================
\bloqueNormativo
%==============================================================================

\subsection{La delimitación legal de la Empresa Deudora}

Análisis exegético de la exigencia de haber sido contribuyente de primera categoría
dentro de los veinticuatro meses anteriores al inicio del procedimiento concursal,
conforme al \art{2}.\verificar{Numeral exacto del art. 2 que define Empresa Deudora en
el texto vigente tras la Ley 21.563.}

\subsection{La noción de Persona Deudora tras la reforma}

Estudio de la inclusión explícita de quienes emiten boletas de honorarios
---contribuyentes de segunda categoría del artículo 42 \Nro 2 de la Ley sobre Impuesto
a la Renta\index[leyes]{Ley sobre Impuesto a la Renta!art. 42 N.º 2}--- dentro del
procedimiento administrativo de renegociación, corrigiendo una exclusión histórica del
foro.\verificar{Referencia precisa al artículo de la Ley sobre Impuesto a la Renta y a
la norma de la Ley 21.563 que incorpora la modificación.}

\subsection{La armonización con los límites de tamaño de la Ley \Nro 20.416}
\pendiente{Desarrollar la remisión y sus problemas de cómputo.}

%==============================================================================
\bloquePractico
%==============================================================================

\subsection{Tabla de calificación del deudor y determinación del procedimiento idóneo}

La siguiente matriz permite clasificar de inmediato al deudor según sus ingresos y
determinar el procedimiento aplicable en la práctica procesal.

\begin{table}[htbp]
\centering
\small
\caption{Calificación del deudor y procedimientos aplicables}
\label{tab:cap02-calificacion}
\begin{tabularx}{\textwidth}{@{}>{\raggedright\arraybackslash}p{0.20\textwidth}
                                 >{\raggedright\arraybackslash}p{0.26\textwidth}
                                 >{\raggedright\arraybackslash}X
                                 >{\raggedright\arraybackslash}X@{}}
\toprule
\textbf{Categoría de deudor} & \textbf{Criterio de ventas anuales (Ley \Nro 20.416)}
& \textbf{Reorganización aplicable} & \textbf{Liquidación aplicable} \\
\midrule
Microempresa    & Hasta \uf{2.400}                    & Simplificada & Simplificada \\
Pequeña empresa & Más de \uf{2.400} y hasta \uf{25.000}  & Simplificada & Simplificada \\
Mediana empresa & Más de \uf{25.000} y hasta \uf{100.000} & Ordinaria    & Ordinaria \\
Gran empresa    & Más de \uf{100.000}                  & Ordinaria    & Ordinaria \\
\addlinespace
Persona Deudora & Sin ingresos de primera categoría   & Renegociación administrativa ante la \superir{} & Simplificada \\
\bottomrule
\end{tabularx}
\end{table}

\begin{alerta}[Cómputo del umbral]
El umbral se mide sobre ventas anuales, no sobre patrimonio ni sobre número de
trabajadores. Antes de invocar la vía simplificada debe acreditarse el nivel de ventas
con la carpeta tributaria del \sii{}, y no con la declaración del cliente.
\end{alerta}

\subsection{Protocolo de acreditación de la calidad del deudor}
\pendiente{Desarrollar el listado documental y los errores frecuentes de calificación.}

%==============================================================================
\bloqueJurisprudencial
%==============================================================================

\subsection{Vigencia temporal de los ingresos de primera categoría}
\pendiente{Sistematizar la jurisprudencia sobre determinación de competencia y tipo de
procedimiento según la vigencia de los ingresos de primera categoría.}

\subsection{Inicio de actividades sin movimiento tributario real}

Doctrina de la \ca{Santiago} conforme a la cual la sola posesión de un inicio de
actividades en primera categoría, sin movimiento tributario real, no priva a una persona
natural de su calidad de \personadeudora{} si se demuestra que su sustento económico
proviene exclusivamente de un contrato de trabajo dependiente o de pensiones de
jubilación.\verificar{Individualizar rol, sala y fecha de los fallos que sostienen esta
doctrina.}
